\PassOptionsToPackage{table,svgnames,dvipsnames}{xcolor}

\usepackage[a-2u]{pdfx} % generate PDF/A: archival compliant, self-contained pdf
\usepackage[utf8]{inputenc}
\usepackage[T1]{fontenc}
\usepackage{amsmath}
\usepackage{amssymb}
\usepackage[sc]{mathpazo}
\usepackage[ngerman,american]{babel}
\usepackage[autostyle]{csquotes}
\usepackage[%
  backend=biber,
  url=false,
  style=alphabetic,
  maxnames=4,
  minnames=3,
  maxbibnames=99,
  giveninits,
  uniquename=init]{biblatex} % TODO: adapt citation style
\usepackage{graphicx}
\usepackage{scrhack} % necessary for listings package
\usepackage{listings}
\usepackage{lstautogobble}
\usepackage{tikz}
\usepackage{pgfplots}
\usepackage{pgfplotstable}
\usepackage{booktabs}
\usepackage[final]{microtype}
\usepackage{caption}
\usepackage{subcaption}
\usepackage[printonlyused]{acronym}
\usepackage{ifthen}
\usepackage{lipsum}
\usepackage{enumitem}
\usepackage{todonotes}
\usepackage{siunitx}
\usepackage{tabularx}
\usepackage{tcolorbox}
\usepackage[table]{xcolor}
\usepackage[toc, nonumberlist]{glossaries}

\setglossarystyle{altlist}
\makeglossaries

\usetikzlibrary{positioning, external, arrows.meta, shapes.geometric, calc, decorations.markings}
% \usetikzlibrary{calc}
% \usetikzlibrary{arrows.meta}
% \usetikzlibrary{decorations.markings}
% Enable externalization
\tikzexternalize[prefix=build/] % Saves the cached figures into the build/ directory

% FIX 1: Remove -output-directory=build. 
% The prefix=build/ already sets the jobname to build/main-figureX, so the output goes to the right place.
\tikzset{external/system call={pdflatex \tikzexternalcheckshellescape -halt-on-error -interaction=batchmode -jobname "\image" "\texsource"}}

% FIX 2: Make todonotes and tikzexternalize play nicely
\makeatletter
\renewcommand{\todo}[2][]{\tikzexternaldisable\@todo[#1]{#2}\tikzexternalenable}
\makeatother

\newcolumntype{L}{>{\raggedright\arraybackslash}X}

\DeclareSourcemap{
  \maps[datatype=bibtex]{
    \map{
       \step[fieldsource=entrykey, match={dealiilibrarysolvergmres}, final]
        \step[fieldset=shorthand, fieldvalue={DII:GMRES}]
    }
  }
}

\hypersetup{hidelinks} % removes colored boxes around references and links

% for fachschaft_print.pdf
\makeatletter
\if@twoside
	\typeout{TUM-Dev LaTeX-Thesis-Template: twoside}
\else
	\typeout{TUM-Dev LaTeX-Thesis-Template: oneside}
\fi
\makeatother

\addto\extrasamerican{
	\def\lstnumberautorefname{Line}
	\def\chapterautorefname{Chapter}
	\def\sectionautorefname{Section}
	\def\subsectionautorefname{Subsection}
	\def\subsubsectionautorefname{Subsubsection}
}

\addto\extrasngerman{
	\def\lstnumberautorefname{Zeile}
}


\providecommand\comment[1]\null

% Themes
\ifthenelse{\equal{\detokenize{dark}}{\jobname}}{%
  % Dark theme
  \newcommand{\bg}{black} % background
  \newcommand{\fg}{white} % foreground
  \usepackage[pagecolor=\bg]{pagecolor}
  \color{\fg}
}{%
  % Light theme
  \newcommand{\bg}{white} % background
  \newcommand{\fg}{black} % foreground
}

\bibliography{bibliography}

\setkomafont{disposition}{\normalfont\bfseries} % use serif font for headings
\linespread{1.05} % adjust line spread for mathpazo font

% Add table of contents to PDF bookmarks
\BeforeTOCHead[toc]{{\cleardoublepage\pdfbookmark[0]{\contentsname}{toc}}}

% Define TUM corporate design colors
% Taken from http://portal.mytum.de/corporatedesign/index_print/vorlagen/index_farben
\definecolor{TUMBlue}{HTML}{0065BD}
\definecolor{TUMSecondaryBlue}{HTML}{005293}
\definecolor{TUMSecondaryBlue2}{HTML}{003359}
\definecolor{TUMBlack}{HTML}{000000}
\definecolor{TUMWhite}{HTML}{FFFFFF}
\definecolor{TUMDarkGray}{HTML}{333333}
\definecolor{TUMGray}{HTML}{808080}
\definecolor{TUMLightGray}{HTML}{CCCCC6}
\definecolor{TUMAccentGray}{HTML}{DAD7CB}
\definecolor{TUMAccentOrange}{HTML}{E37222}
\definecolor{TUMAccentGreen}{HTML}{A2AD00}
\definecolor{TUMAccentLightBlue}{HTML}{98C6EA}
\definecolor{TUMAccentBlue}{HTML}{64A0C8}

% Settings for pgfplots
\pgfplotsset{compat=newest}
\pgfplotsset{
  % For available color names, see http://www.latextemplates.com/svgnames-colors
  cycle list={TUMBlue\\TUMAccentOrange\\TUMAccentGreen\\TUMSecondaryBlue2\\TUMDarkGray\\},
}

% Settings for lstlistings
\lstset{%
  basicstyle=\ttfamily,
  columns=fullflexible,
  autogobble,
  keywordstyle=\bfseries\color{TUMBlue},
  stringstyle=\color{TUMAccentGreen},
  captionpos=b
}


% Define a custom environment called 'insetBox'
\newtcolorbox{insetBox}[1][]{
  shield externalize,
  colback=TUMLightGray!20,       % Light gray background (10% gray, 90% white)
  colframe=TUMBlue,      % Slightly darker gray border
  title=Performance Tip, % Default title for the box
  fonttitle=\bfseries,   % Bold title font
  coltitle=black,        % Black text for the title
  arc=3pt,               % Gently rounded corners
  boxrule=0.5pt,         % Thin, unobtrusive border
  #1                     % Allows you to pass extra options locally if needed
}